\documentclass[12pt,letterpaper]{article}
\usepackage[OT1]{fontenc}
\renewcommand{\rmdefault}{cmr}
\renewcommand{\sfdefault}{cmss}
\renewcommand{\ttdefault}{cmtt}
\usepackage{amsmath,amssymb}
\usepackage[left=1.25in,right=1.25in,top=1.05in,bottom=1.38in,footskip=0.40in]{geometry}
\usepackage{enumitem,graphicx,needspace}
\setlength{\parindent}{0pt}
\setlength{\parskip}{0.45em}
\setlist{topsep=0.35em,itemsep=0.3em,parsep=0pt}
\setlength{\emergencystretch}{2em}
\pagestyle{plain}
\newcommand{\problemheading}[2][5\baselineskip]{\Needspace{#1}\par\noindent\hspace*{-0.625in}\textbf{#2}\par\nobreak\vspace{0.3em}}
\newcommand{\examspace}[1]{\par\vspace*{#1cm}}
\begin{document}
\thispagestyle{empty}
\begin{center}
\textbf{COMP 2711}\\
\textbf{Discrete Math Tools for Computer Science}\\
{\large\bfseries Mock Midterm Examination 08}\\[1em]
\textbf{Practice material -- not an official examination}\\[0.7em]
Time allowed: 120 minutes \hfill Total: 100 points
\end{center}
\vspace{1.1cm}
\begin{center}\fbox{\fbox{\begin{minipage}{0.93\textwidth}
Name: \hrulefill\quad Student ID: \hrulefill\\[0.6em]
Email: \hrulefill\quad Lecture: \hrulefill
\end{minipage}}}\end{center}
\vspace{0.6cm}
\textbf{Instructions}
\begin{itemize}[leftmargin=1.3em]
\item This closed-book practice exam consists of 10 problems. You may use one A4 single-sided handwritten cheat sheet.
\item Answer all problems in the spaces provided. You may also use the backs of pages and the draft pages. Each problem starts on a new page; some answers require only a few lines.
\item Give a proof, justification, or counterexample when requested. State any additional notation you introduce.
\item Answers may be left in exact form using binomial coefficients, factorials, and $C(n,k)$ or $P(n,k)$. For example, $\binom{5}{3}+\binom{4}{2}$ is acceptable. Avoid nonstandard notation such as $nCk$ or $nPk$.
\item The separate solutions document contains derivations and marking notes. Attempt the exam before consulting it.
\end{itemize}
\vspace{1.0cm}
\begin{center}
\Large
\setlength{\tabcolsep}{2pt}
\begin{tabular*}{\textwidth}{@{\extracolsep{\fill}}|l|*{11}{c|}}\hline
Problems & 1 & 2 & 3 & 4 & 5 & 6 & 7 & 8 & 9 & 10 & Total\\\hline
Points & 10 & 13 & 12 & 16 & 10 & 9 & 7 & 6 & 7 & 10 & 100\\\hline
Score & & & & & & & & & & & \\[0.7em]\hline
\end{tabular*}
\end{center}
\vspace{0.8cm}

\clearpage
\vspace*{0.21in}
\problemheading{Problem 1: Propositional logic and proofs [10 points]}
\begin{list}{}{\setlength{\leftmargin}{0.25in}\setlength{\rightmargin}{0.25in}}\item[]
\begin{description}[font=\normalfont,leftmargin=1.2em,labelsep=0.4em]
\item[(a)] \textbf{[3 points]} Using propositional laws, prove $(p\leftrightarrow q)\land p\equiv p\land q$.
\examspace{4}
\Needspace{\dimexpr 3cm+4\baselineskip\relax}
\item[(b)] \textbf{[2 points]} Classify each as a proposition or not: (i) $2+3=5$; (ii) $x+3=5$ with unassigned $x$; (iii) \emph{Close the door}; (iv) \emph{Is the door closed?} Give the truth value whenever meaningful. Each is worth $1/2$ point.
\examspace{3}
\Needspace{\dimexpr 6cm+4\baselineskip\relax}
\item[(c)] \textbf{[5 points]} For an integer $n$, prove by a cycle of implications that the following are equivalent: (I) $n$ is odd; (II) $n+1$ is even; (III) $n^2$ is odd.
\examspace{6}
\end{description}
\end{list}
\clearpage
\vspace*{0.21in}
\problemheading{Problem 2: Predicate logic [13 points]}
\begin{list}{}{\setlength{\leftmargin}{0.25in}\setlength{\rightmargin}{0.25in}}\item[]
\begin{description}[font=\normalfont,leftmargin=1.2em,labelsep=0.4em]
\item[(a)] \textbf{[4 points]} Variables range over $\mathbb R$. Translate (i) every nonzero real has a real multiplicative inverse; (ii) some negative real has a square equal to itself. Use unrestricted quantifiers. Then explain why $\exists x(x<0\to x^2=x)$ is an incorrect translation of (ii).
\examspace{4}
\Needspace{\dimexpr 4cm+4\baselineskip\relax}
\item[(b)] \textbf{[4 points]} Negate $\forall x\exists y(P(x,y)\land\forall z(Q(y,z)\to R(x,z)))$, placing negations immediately before predicates only.
\examspace{4}
\Needspace{\dimexpr 6cm+4\baselineskip\relax}
\item[(c)] \textbf{[5 points]} Explain why instantiating $\exists xP(x)$ to $P(a)$ and then universally generalizing to $\forall xP(x)$ is invalid, using a countermodel (2 points). Give a valid formal derivation of $\forall x(P(x)\lor Q(x))$ from $\forall xP(x)$ (3 points).
\examspace{6}
\end{description}
\end{list}
\clearpage
\vspace*{0.21in}
\problemheading{Problem 3: Sets and functions [12 points]}
\begin{list}{}{\setlength{\leftmargin}{0.25in}\setlength{\rightmargin}{0.25in}}\item[]
\begin{description}[font=\normalfont,leftmargin=1.2em,labelsep=0.4em]
\item[(a)] \textbf{[2 points]} List $\mathcal P(\varnothing)$ and $\mathcal P(\mathcal P(\varnothing))$, keeping all braces explicit.
\examspace{3}
\Needspace{\dimexpr 5cm+4\baselineskip\relax}
\item[(b)] \textbf{[3 points]} Prove $(0,1)$ is uncountable by diagonalizing against an alleged list of its canonical decimal expansions, chosen not to end in repeating 9s. Use only digits 4 and 5 in the new number and address decimal ambiguity.
\examspace{5}
\Needspace{\dimexpr 1cm+4\baselineskip\relax}
\item[(c)] \textbf{[7 points]} Write T or F for each statement. No justification is required. Each item is worth 1 point.\begin{enumerate}[label=(\roman*)]\item An axiom is a statement assumed to be true.\item A theorem is a statement established by proof.\item A lemma is a theorem used to help prove other theorems.\item A corollary follows directly from a previously proved theorem.\item Every conjecture already has a proof.\item A proof is a valid argument establishing a theorem.\item The self-referential sentence \emph{This sentence is false} is a proposition in the classical definition used here.\end{enumerate}
\examspace{1}
\end{description}
\end{list}
\clearpage
\vspace*{0.21in}
\problemheading{Problem 4: Short counting calculations [16 points]}
\begin{list}{}{\setlength{\leftmargin}{0.25in}\setlength{\rightmargin}{0.25in}}\item[]
Give exact answers. Factorials, powers, and binomial coefficients may be left unevaluated. No explanation is required unless requested.
\begin{description}[font=\normalfont,leftmargin=1.2em,labelsep=0.4em]
\item[(a)] \textbf{[2 points]} Tickets have format LDD or LLD, with letters from $\{A,B,C\}$ and digits from $\{0,1,2,3\}$. Within a ticket, letters cannot repeat and digits cannot repeat. Count the tickets.
\examspace{1.5}
\Needspace{\dimexpr 1.5cm+4\baselineskip\relax}
\item[(b)] \textbf{[2 points]} Choose and order four distinct people among eight, requiring $A$ to occur and $B$ not to occur. Count the ordered selections.
\examspace{1.5}
\Needspace{\dimexpr 1.5cm+4\baselineskip\relax}
\item[(c)] \textbf{[2 points]} Seven distinct people stand in a line with $A$ and $B$ not adjacent. Count the lines.
\examspace{1.5}
\Needspace{\dimexpr 1.5cm+4\baselineskip\relax}
\item[(d)] \textbf{[2 points]} Six distinct people sit at equally spaced places around a circular table, with $A$ and $B$ opposite. Rotations are equivalent; reflections are different. Count the seatings.
\examspace{1.5}
\Needspace{\dimexpr 1.5cm+4\baselineskip\relax}
\item[(e)] \textbf{[2 points]} Count length-nine bit strings containing four ones and ending in 00.
\examspace{1.5}
\Needspace{\dimexpr 1.5cm+4\baselineskip\relax}
\item[(f)] \textbf{[2 points]} An unlimited fruit supply has six types, $A$ through $F$. Choose ten pieces, with at least two of type $A$ and three of type $B$. The other four types may be absent. Fruit of a type is identical and order does not matter. Count the choices.
\examspace{1.5}
\Needspace{\dimexpr 1.5cm+4\baselineskip\relax}
\item[(g)] \textbf{[2 points]} Find the coefficient of $x^2y^4$ in $(3x-y)^6$.
\examspace{1.5}
\Needspace{\dimexpr 1.5cm+4\baselineskip\relax}
\item[(h)] \textbf{[2 points]} Count right/up unit-step lattice paths from $(0,0)$ to $(4,4)$ that pass through $(2,1)$.
\examspace{1.5}
\end{description}
\end{list}
\clearpage
\vspace*{0.21in}
\problemheading{Problem 5: Combinatorial proof [10 points]}
\begin{list}{}{\setlength{\leftmargin}{0.25in}\setlength{\rightmargin}{0.25in}}\item[]
\begin{description}[font=\normalfont,leftmargin=1.2em,labelsep=0.4em]
\item[(a)] \textbf{[7 points]} Use the choices made when multiplying $n$ binomial factors to prove, for $n\ge0$,\[(x+y)^n=\sum_{k=0}^{n}\binom nk x^{n-k}y^k.\]
\examspace{5}
\Needspace{\dimexpr 3cm+4\baselineskip\relax}
\item[(b)] \textbf{[3 points]} For $0\le r\le n$, deduce the coefficient of $x^{2n-2r}y^r$ in $(x^2-3y)^n$.
\examspace{3}
\end{description}
\end{list}
\clearpage
\vspace*{0.21in}
\problemheading{Problem 6: Inclusion-exclusion [9 points]}
\begin{list}{}{\setlength{\leftmargin}{0.25in}\setlength{\rightmargin}{0.25in}}\item[]
\begin{description}[font=\normalfont,leftmargin=1.2em,labelsep=0.4em]
\item[] \textbf{[9 points]} A nine-element set contains three specified disjoint pairs. Use inclusion--exclusion to count its five-element subsets in which no specified pair has both members selected.
\examspace{7}
\end{description}
\end{list}
\clearpage
\vspace*{0.21in}
\problemheading{Problem 7: Pigeonhole principle [7 points]}
\begin{list}{}{\setlength{\leftmargin}{0.25in}\setlength{\rightmargin}{0.25in}}\item[]
\begin{description}[font=\normalfont,leftmargin=1.2em,labelsep=0.4em]
\item[] \textbf{[7 points]} A collection contains ten red, eight blue and six green tokens. Tokens are selected without replacement. Find the smallest number that guarantees at least six red, or at least four blue, or at least three green, and prove that your number is minimal.
\examspace{7}
\end{description}
\end{list}
\clearpage
\vspace*{0.21in}
\problemheading{Problem 8: Unequal success probabilities [6 points]}
\begin{list}{}{\setlength{\leftmargin}{0.25in}\setlength{\rightmargin}{0.25in}}\item[]
Three sensors independently pass a check with probabilities $1/2,1/3,1/4$, respectively. All sensors are checked once; outcomes are the three-symbol pass/fail strings with product masses. Let $X$ be the number that pass.
\begin{description}[font=\normalfont,leftmargin=1.2em,labelsep=0.4em]
\item[(a)] \textbf{[1 points]} Find $P(X=3)$.
\examspace{1.2}
\Needspace{\dimexpr 1.2cm+4\baselineskip\relax}
\item[(b)] \textbf{[1 points]} Find $P(X=0)$.
\examspace{1.2}
\Needspace{\dimexpr 1.2cm+4\baselineskip\relax}
\item[(c)] \textbf{[1 points]} Find $P(X=1)$.
\examspace{1.2}
\Needspace{\dimexpr 1.2cm+4\baselineskip\relax}
\item[(d)] \textbf{[1 points]} Find $P(X\ge1)$.
\examspace{1.2}
\Needspace{\dimexpr 2.0cm+4\baselineskip\relax}
\item[(e)] \textbf{[2 points]} Is the taught binomial formula justified by this trial model? State both model requirements and identify which one fails.
\examspace{2.0}
\end{description}
\end{list}
\clearpage
\vspace*{0.21in}
\problemheading{Problem 9: Message evidence [7 points]}
\begin{list}{}{\setlength{\leftmargin}{0.25in}\setlength{\rightmargin}{0.25in}}\item[]
An archive contains 360 spam and 240 legitimate messages. A word $W$ appears in 90 spam and 12 legitimate messages. Choose a message uniformly. A second marker $V$ has conditional probabilities $2/5$ in spam and $1/6$ in legitimate messages; $W,V$ are independent conditional on each message type.
\begin{description}[font=\normalfont,leftmargin=1.2em,labelsep=0.4em]
\item[(a)] \textbf{[3 points]} Find $P(W)$ and $P(\text{spam}\mid W)$.
\examspace{3.5}
\Needspace{\dimexpr 4.5cm+4\baselineskip\relax}
\item[(b)] \textbf{[4 points]} Given that both markers occur, find the spam posterior. Explain which independence permits multiplying the marker likelihoods.
\examspace{4.5}
\end{description}
\end{list}
\clearpage
\vspace*{0.21in}
\problemheading{Problem 10: Dependent counts and products [10 points]}
\begin{list}{}{\setlength{\leftmargin}{0.25in}\setlength{\rightmargin}{0.25in}}\item[]
A coin is tossed $n\ge1$ times independently, with heads probability $2/3$. Let $H,T$ count heads and tails. For part (d), $H^{\prime}$ counts heads in a separate independent batch of $n$ tosses with the same law.
\begin{description}[font=\normalfont,leftmargin=1.2em,labelsep=0.4em]
\item[(a)] \textbf{[2 points]} Find $E(H)$ and $E(T)$.
\examspace{2.0}
\Needspace{\dimexpr 2.0cm+4\baselineskip\relax}
\item[(b)] \textbf{[2 points]} Find $E(HT)$ using $T=n-H$.
\examspace{2.0}
\Needspace{\dimexpr 2.0cm+4\baselineskip\relax}
\item[(c)] \textbf{[2 points]} Find $E(H+T)$ and $V(H+T)$.
\examspace{2.0}
\Needspace{\dimexpr 4.5cm+4\baselineskip\relax}
\item[(d)] \textbf{[4 points]} Show that $H,T$ are dependent. Find $E(HH^{\prime})$, and justify product-expectation factorization from the joint-mass definition for finite independent variables.
\examspace{4.5}
\end{description}
\end{list}
\clearpage\textbf{Draft Paper}\par\textit{Mock examination working space.}\vfill
\clearpage\textbf{Draft Paper}\par\textit{Mock examination working space.}\vfill
\end{document}
